\section{对数几何导引（Tingting Fang 与 Zijian Han）}%
\label{sec:Intro to log geometry (Tingting Fang and Zijian Han)}

\subsection{动机}%
\label{sub:Motivation}

在计数几何中，我们将考虑紧化模空间上的相交数。这将要求我们的对象以某种方式退化，例如定义域曲线变为结点曲线、映射落入边界除子，或者使目标退化。在普通几何中，我们把这些情形视为奇异的，但在对数几何中，有一种办法把这些情形看作\textit{对数光滑}的。在对数几何中，我们将用函数连同函数在边界上如何为零的信息（也就是态射 \(\alpha \colon M_X \to \mc{O}_X\)）来替代函数。

\begin{exm}
    考虑 \(\A^1\) 上的曲线族 \((xy=t)\)，其中心纤维是结点的。在对数几何中，如果我们在 \(\A^1\) 上赋予除子对数结构，那么这个族是对数光滑的。
\end{exm}

\subsection{幺半群}%
\label{sub:Monoids}

\begin{defn}
    \textit{幺半群}是带单位元的交换半群。
\end{defn}

对于一个幺半群 \(P\)，我们可以通过标准的 Grothendieck 构造关联一个群 \(P^{\mr{gp}}\)：
\[ P^{\mr{gp}} = \ab\{(a,b) \mid (a,b) \sim (c,d) \text{ if there exists } s \in P \text{ such that } a+d+s = b+c+s \}. \]
\begin{defn}
如果自然态射 \(P \hookrightarrow P^{\mr{gp}}\) 是单射，我们称 \(P\) 是\textit{整的}。
\end{defn}

\begin{exm}
    如果 \(P\) 有生成元 \(a,b,c\) 与关系 \(a+c=b+c\)，那么在 \(P^{\on{gp}}\) 中我们有 \(a = b\)，故 \(P\) 不是整的。
\end{exm}

\begin{defn}
    幺半群 \(P\) 称为
    \begin{itemize}
        \item \textit{细的}，如果 \(P\) 是有限生成且整的；
        \item \textit{饱和的}，如果对所有 \(p \in P^{\mr{gp}}\) 与所有自然数 \(n\)，由 \(np \in P\) 可得 \(p \in P\)；
        \item \textit{无挠的}，如果 \(P^{\mr{gp}}\) 是无挠的；
        \item \textit{尖锐的}，如果单位元是唯一的可逆元。
    \end{itemize}
\end{defn}

\subsection{对数结构}%
\label{sub:Log structures}

设 \(X\) 为概形。
\begin{defn}
    \(X\) 上的\textit{预对数结构}是 \(X\) 的小平展站点上的一个幺半群层 \(M_X\) 以及一个态射 \(\alpha \colon M_X \to \mc{O}_X\)。如果 \(\alpha^{-1}(\mc{O}_X^{\times}) \simeq \mc{O}_X^{\times}\)（其中同构由 \(\alpha\) 给出），则称其为\textit{对数结构}。
\end{defn}

现在考虑正合序列
\[ 0 \to \alpha^{-1}(\mc{O}_X^{\times}) \to M_X \to \bar{M}_X \to 0, \]
这里的商称为\textit{特征层}或\textit{鬼层}。

\begin{exm}
    令 \(M_X = \mc{O}_X^{\times}\) 且 \(\alpha\) 为恒等映射。此时对数结构不提供任何额外的信息。
\end{exm}

\begin{exm}
    设 \(X\) 为正则簇，\(D\) 为闭子概形。令 \(U = X \setminus D\)。定义
    \[ M_X = \mc{O}_X \cap j_* \mc{O}_U^{\times}, \]
    其中 \(j \colon U \to X\) 为开浸入。如果 \(D\) 是一个除子，则称之为\textit{除子对数结构}。

    例如，如果 \(X = \A^1\) 且 \(D = \ab\{0\}\)，那么特征层除原点外处处为零。这里我们计算
    \[ M_X = \ab\{ \varphi \cdot t^n \mid \varphi \in \mc{O}_X^{\times} \}. \]
    其茎除原点外处处为 \(\mc{O}_X^{\times}\)，而在原点处为 \(\mc{O}_{\A^1, 0}^{\times} \oplus \N\)。
\end{exm}

\begin{exm}
    现在令 \(X = \A^2\)，并令 \(D\) 为坐标轴之并。则
    \begin{align*}
        M_{X,D} &= \ab\{f \in k[x,y] \mid f \text{ is invertible outside } D \} \\
        &= \ab\{\varphi \cdot x^a \cdot y^b\}.
    \end{align*}
    注意到在一般点处，我们看到 \(M_p = \mc{O}_{X,p}^{\times}\)。在原点处，我们看到 \(\mc{O}_{X,p}^{\times}\oplus \N^2\)。在 \(x\)-轴的一般点处，\(x^a\) 可逆，故茎为 \(\mc{O}_{X,p}^{\times}\oplus \N \cdot b\)。类似地，我们可以计算 \(y\)-轴上一般点处的茎。
\end{exm}

正如我们可以给预层关联一个层一样，我们也会给预对数结构关联一个对数结构。首先，范畴 \(\mc{C}\) 中的推出记作 \(B \sqcup_A C\)，它是图表
\begin{equation*}
\begin{tikzcd}
    A \ar{r}{f} \ar{d}[swap]{g} & B \\
    C.
\end{tikzcd}
\end{equation*}
的余极限。当 \(\mc{C}\) 是交换幺半群范畴时，
\[ B \sqcup_A C = (B \oplus C)/\sim, \qquad (f(a) + b, c) \sim (b, g(a)+c). \]

\begin{exm}
    考虑态射 \(2 \colon \N \to \N\) 与 \(3 \colon \N \to \N\)。此时推出是幺半群
    \[ P = \frac{\N \oplus \N}{(3, 0) \sim (0,2)} \simeq \N \setminus \{1\}. \]
\end{exm}

如果把 \(\mc{C}\) 换成交换幺半群层的范畴，我们可以先在截面上作推出，然后再层化。

\begin{defn}
    预对数结构 \(\alpha \colon M_X \to \mc{O}_X\) 的\textit{相伴对数结构}是推出
    \[ M_X \sqcup_{\alpha^{-1}(\mc{O}_X^{\times})} \mc{O}_X^{\times}. \]
    到 \(\mc{O}_X\) 的映射由 \((a,b) \to \alpha(a) b\) 给出。
\end{defn}

\begin{defn}
    对数概形的态射 \((X, M_X) \to (Y, M_Y)\) 由以下组成：
    \begin{enumerate}
        \item 底层概形的态射 \(\ul{f} \colon \ul{X} \to \ul{Y}\)；
        \item 映射 \(f^{\flat} \colon f^* M_Y \to M_X\)，其中 \(f^* M_Y\) 是与 \(f^{-1} M_Y\) 相伴的对数结构。
    \end{enumerate}
\end{defn}

\begin{defn}
    对数概形 \((X, \alpha \colon M_X \to \mc{O}_X)\) 的\textit{图卡}是从常值幺半群出发的态射 \(P_X \to M_X\)，使得 \(M_X \simeq P_X^a \simeq P_X \oplus_{\alpha_P^{-1}(\mc{O}_X^{\times})} \mc{O}_X^{\times}\)，其中 \(\alpha_P \colon P_X \to M_X \to \mc{O}_X\)。
\end{defn}

图卡在指定对数结构时很有用。例如，如果我们有 \(P_X \to \mc{O}_X\)，那么 \(P_X^a\) 给出 \(X\) 上的一个对数结构。

\begin{exm}
    平凡对数结构由图卡 \(P = \ab\{0\}\) 给出。
\end{exm}

\begin{exm}
    令 \(X = \Spec k\)，并由 \(\alpha_{\N} \colon \N \to k\) 给出对数结构，它把 \(0\) 送到 \(1\)、把 \(1\) 送到 \(0\)。这称为\textit{标准对数点}，也可以通过把 \(\A^1\) 上的除子对数结构限制到原点而得到。
\end{exm}

\begin{exm}
    考虑把 \(n \mapsto t^n\) 的图卡 \(\N \to \mc{O}_{\A^1}\)。它恢复了 \(\A^1\) 上相对于原点的除子对数结构。
\end{exm}

\begin{exm}
    考虑把 \((a,b) \mapsto x^a y^b\) 的图卡 \(\N^2 \to \mc{O}_{\A^2}\)。它恢复了 \(\A^2\) 上相对于环簇边界的除子对数结构。
\end{exm}

\begin{defn}
    对数映射 \(f \colon (X, M_X) \to (Y, M_Y)\) 的\textit{图卡}是一个三元组 \((P_X\to M_X, Q_Y \to M_Y, Q \to P)\)，其中 \(P_X \to M_X\) 与 \(Q_Y \to M_Y\) 是图卡，且图表
    \begin{equation*}
    \begin{tikzcd}
        Q_X \ar{r} \ar{d} & P_X \ar{d} \\
        f^* M_Y \ar{r} & M_X
    \end{tikzcd}
    \end{equation*}
    交换。
\end{defn}

\begin{defn}
    称对数概形 \((X, M_X)\) 是\textit{有限生成的、细的、整的、饱和的}，如果存在一个覆盖 \(U \to X\) 以及平展局部图卡 \(P_U \to M_U\)，使得 \(P\) 分别是有限生成的、细的、整的或饱和的。
\end{defn}

\subsection{对数光滑性}%
\label{sub:Log smoothness}

\begin{defn}
    考虑交换图表
    \begin{equation*}
    \begin{tikzcd}
        T_0 \ar{r}{\phi} \ar{d}{j} & X \ar{d}{f} \\
        T_1 \ar{r}{\psi} & Y,
    \end{tikzcd}
    \end{equation*}
    其中 \(j\) 是严格闭浸入（这里指的是 \(M_{T_0} = j^* M_{T_1}\)）。称细对数概形的态射 \(f \colon X \to Y\) 是\textit{对数光滑（相应地，对数平展）}的，如果
    \begin{enumerate}
        \item 底层映射 \(\ul{f} \colon \ul{X} \to \ul{Y}\) 是局部有限展示的；
        \item 对任意如上交换图表，在 \(T_1\) 上平展局部地存在（相应地，存在唯一）态射 \(g \colon T_1 \to X\)，使得 \(\phi = g \circ j\) 且 \(\psi = f \circ g\)。
    \end{enumerate}
\end{defn}

现在我们将给出对数光滑性的一个图卡局部判据。

\begin{thm}[{\cite[Theorem~3.5]{Kato1989}}]
    设 \(f \colon X \to Y\) 为细对数概形的态射。假设我们有图卡 \(Q \to M_Y\)，其中 \(Q\) 是细幺半群。则 \(f\) 对数光滑当且仅当在 \(X\) 上平展局部地存在 \(f\) 的一个延拓 \(Q \to M_Y\) 的图卡，满足
    \begin{enumerate}
        \item \(Q^{\mr{gp}} \to P^{\mr{gp}}\) 的核与余核的挠部分是有限群，且其阶在 \(X\) 上可逆；
        \item 诱导态射 \(\ul{X} \to \ul{Y} \times_{\Spec \Z[Q]} \Spec \Z[P]\) 在经典意义下是平展的。
    \end{enumerate}
\end{thm}

\begin{exm}
    设 \(X\) 是由幺半群 \(P\) 给出的仿射环簇 \(\Spec \C[P]\)。注意 \(P^a \simeq M_{X, X \setminus T}\) 诱导出除子对数结构。然后考虑映射 \(f \colon X \to \mr{pt}\)，其中该点带有平凡对数结构。我们注意到态射 \(0 \to P^{\mr{gp}}\) 没有核且余核无挠，并且映射
    \[ \ul{X} \to \Spec \C \times_{\Spec \Z} \Spec \Z[P] = \Spec \C[P] \]
    是同构，因此任何环簇在带平凡对数结构的点上都是对数光滑的。
\end{exm}

从这个例子来看，我们在直观上可以把对数光滑性看作相对于基底是局部环簇的。

\begin{exm}
    回忆曲线族 \(xy = t\)，其对数结构相对于中心纤维，底空间 \(\A^1\) 带有相对于 \(0\) 的除子对数结构。我们将考虑的图卡映射实际上是 \(\Delta \colon \N \to \N^2\)。取群化之后，映射变为 \(\Delta \colon \Z \to \Z^2\)。其核是平凡的且余核无挠，于是我们验证了第一个条件。然后我们考虑映射
    \[ \ul{X} \to \Spec \C[t] \times_{\Spec \Z[\N]} \Spec \Z[\N^2] = \Spec \frac{\C[u,v,t]}{uv-t}, \]
    它是同构，故此族是对数光滑的。
\end{exm}

\subsection{对数结构之栈}%
\label{sub:Stack of log structures}

对数结构之栈及其与对数光滑性的关系归功于 Olsson~\cite{Olsson2003}。

设 \(S\) 为细对数概形。我们将考虑范畴 \(\ms{Log}_S\)，其对象是细对数概形 \((T, M_T) \to S\)，其态射是 \(S\) 上的严格态射 \((T, M_T) \to (T', M_{T'})\)。
\begin{enumerate}
    \item 范畴 \(\ms{Log}_S\) 是 \(\ms{Sch}_{/S}\) 上以群胚纤维化的范畴，它分类所有映到 \(S\) 的细对数概形；
    \item \(\ms{Log}_S\) 是 \(S\) 上局部有限型的代数栈；
    \item 态射 \(f \colon X \to S\) 对数光滑当且仅当 \(X \to \ms{Log}_S\) 光滑。把光滑换成平展，结论同样成立。
\end{enumerate}

\subsection{热带化}%
\label{sub:Tropicalization}

设 \(X\) 为细对数概形。则 \(X\) 的热带化定义为余极限
\[\Sigma(X) \coloneqq \bigsqcup_{x \in X} \Hom(\bar{M}_{X,x}, \R_{\geq 0}) \biggm/ \sim, \]
其中等价关系如下给出：若点 \(x\) 特化到 \(y\)，则诱导映射 \(\sigma_x \to \sigma_y\) 给出从 \(\sigma_x\) 到其像的等价。

现在假设 \(X\) 有一个由 \(X_i\) 给出的分层，使得 \(\bar{M}_{|X_i}\) 是常值的。对任意 \(X_i\)，所有点 \(x \in X_i\) 都满足 \(x \in \ol{\ab\{\eta_{X_i}\}}\)，因此我们只需考虑各层片的泛点，然后再粘合各个面。

\begin{exm}
    考虑带平凡对数结构的 \(X\)。注意到对任意点 \(x \in X\)，我们有 \(\bar{M}_{X,x} = 0\)，故 \(\sigma_X = \ab\{0\}\)。于是 \(\Sigma(X) = \ab\{0\}\)。
\end{exm}

\begin{exm}
    考虑相对于原点的 \(\A^1\)（带除子对数结构）。易见 \(\sigma_0 = \Hom(\N, \R_{\geq 0}) = \R_{\geq 0}\) 且 \(\sigma_{\eta} = \ab\{0\}\)，故热带化为 \(\R_{\geq 0}\)。
\end{exm}

\begin{exm}
    考虑带环簇对数结构（相对于 \(xy = 0\) 的除子对数结构）的 \(\A^2\)。则 \(\sigma_{(0,0)} = \R^2_{\geq 0}\)、\(\sigma_{(x,0)} = \R_{\geq 0}\)、\(\sigma_{(0,y)} = \R_{\geq 0}\) 且 \(\sigma_{(x,y)} = \ab\{0\}\)，于是我们得到 \(\A^2\) 作为环簇的扇，并将其视为锥复形。
\end{exm}

\begin{exm}
    设 \(X_{\Sigma}\) 为环簇，带有相对于 \(X_{\Sigma} \setminus T\) 的除子对数结构。则
    \[ \Sigma(X_{\Sigma}) = \Sigma, \]
    于是一个可用的口号是：把环簇热带化可以恢复普通环簇几何的结构（只是去掉了扇在某个实向量空间中的嵌入）。
\end{exm}

\begin{exm}
    设 \(X = \Spec k[x,y,t]/(xy-t^{\ell})\)，带有相对于 \(xy = 0\) 的除子对数结构。则有一个由
    \[ Q = \ab<e_x, e_y, e_t \mid e_x + e_y = \ell e_t>. \]
    给出的图卡。注意 \(M_X = Q^a = Q \oplus_{\alpha^{-1}(\mc{O}_X^{\times})} \mc{O}_X^{\times}\) 且 \(\bar{M}_X = Q\)。因此，在奇点 \(q\) 处我们有
    \begin{align*}
        \sigma_q &= \Hom(Q, \R_{\geq 0}) \\
        &= \ab\{ (\rho_x, \rho_y, \rho_t) \in \R_{\geq 0}^3 \mid \rho_x + \rho_y = \ell \rho_t \} \\
        &= \ab\{ (\rho_t, \rho_x) \in \R_{\geq 0}^2 \mid 0 \leq \rho_x \leq \ell \rho_t \}
    \end{align*}
    这恢复了把 \(X\) 表示为环簇的锥。在 \(x = 0\) 的泛点 \(\eta_x\) 处，我们看到 \(\bar{M}_{X, \eta_x} = \N e_t\)，因此 \(\sigma_{\eta_x} = \R_{\geq 0}\)。类似地，我们看到 \(\sigma_{\eta_y} = \R_{\geq 0}\)。剩下的问题是粘合。

    注意 \(\bar{M}_{X, q} \to \bar{M}_{X, \eta_x}\) 由 \(e_t \mapsto e_t\)、\(e_x \mapsto \ell e_t\) 与 \(e_y \mapsto 0\) 给出，因此映射 \(\sigma_{\eta_x} \to \sigma_q\) 把 \(\rho_t\) 送到 \((\rho_t, \rho_x)\)，其中 \(\rho_x = \ell \rho_t\)。

    态射 \(X \to \A^1_t\) 给出一个中心纤维为结点的曲线族，其热带化由把锥投影到 \(\R_{\geq 0}\) 给出。
\end{exm}


\subsection{对数曲线}%
\label{sub:Log curves}

从现在起，所有幺半群与对数概形都将是细的（甚至是细饱和的）。这有两个原因：
\begin{enumerate}
    \item 如果 \(P\) 是细幺半群，则其对偶是 fs 且尖锐的；
    \item 如果 \(X = (\ul{X}, M_X)\) 是细对数概形，则对任意几何点 \(\bar{x} \to X\)，存在 \(\bar{x}\) 的一个具有图卡的平展邻域。
\end{enumerate}

现在我们将讨论幺半群之间映射的各种几何性质。
\begin{defn}
    称映射 \(\theta \colon P \to Q\) 是\textit{整的}，如果映射 \(\Spec\C[Q] \to \Spec \C[P]\) 是平坦的。
\end{defn}

例如，任何非零映射 \(\N \to \N\) 与对角映射 \(\Delta \colon\N \to \N^2\) 都是整的。如果一个映射是整的，那么对任意 \(P \to P'\)，幺半群 \(P' \oplus_P Q\) 是细的。

\begin{exm}
    推出未必是细的，例如考虑
    \[ \ab< a ,b,c \mid a+c=b+c> = \N^3 \oplus_{\N^2} \N. \]
\end{exm}

\begin{defn}
    称细对数概形之间的映射 \(f \colon X \to Y\) 是整的，如果在 \(X\) 的每个几何点处，映射
    \[ (f^*\bar{M}_Y)_{\bar{x}} \to \bar{M}_{X,\bar{x}} \]
    是整的。
\end{defn}

\begin{defn}
    称 fs 幺半群之间的映射 \(P \to Q\) 是\textit{饱和的}，如果对所有 \(P \to P'\)，幺半群 \(P' \oplus_P Q\) 也是 fs 的。
\end{defn}

\begin{exm}
    映射 \(\Delta \colon \N \to \N^2\) 是饱和的，但 \(\ell \colon \N \to \N\) 不是。
\end{exm}

\begin{defn}
    fs 对数概形 \(S\) 上的\textit{对数曲线}是 fs 对数概形之间的一个对数光滑且整的映射 \(\pi \colon \mc{C} \to S\)，使得其底层几何纤维是既约连通曲线。
\end{defn}

\begin{prop}[{\cite{Kato1989}}]
    设 \(f \colon X \to Y\) 是 fs 对数概形之间对数光滑且整的映射。则 \(\ul{f}\) 是平坦的。
\end{prop}

\begin{thm}[{\cite{FKato2000}}]
    对数曲线的底层族是结点曲线族。
\end{thm}

\begin{exm}
    令 \(P = \N \setminus \ab\{1\}\)。则 \(\Spec (P \to k[P])\) 在 \(k\) 上是对数光滑的，但 \(\Spec k[P]\) 是一个尖点。
\end{exm}

\begin{defn}
    fs 基底 \(S\) 上的\textit{\(n\) 点标记对数曲线}由数据 \((\pi \colon \mc{C} \to S), \vec{p} = \ab\{p_i\}_{i=1}^n\) 组成，使得
    \begin{enumerate}
        \item 底层族是预稳定曲线；
        \item \(\pi\) 是 fs 对数概形之间的对数光滑映射；
        \item 在 \(\ul{\pi}\) 的光滑轨迹 \(U\) 上，我们有
            \[ \bar{M}_{\mc{C}}|_U \simeq \pi^* \bar{M}_S \oplus \bigoplus_{i=1}^n p_i \ul{\N}. \]
    \end{enumerate}
\end{defn}

\subsection{预稳定曲线上的典范对数结构}%
\label{sub:Canonical log structures on prestable curves}

设 \(\mf{M}_{g,n}\) 为 \(n\) 点标记预稳定曲线的模。则其奇异轨迹是正规交叉除子，故有典范对数结构
\[ M_{\mf{M}_{g,n}}^{\on{can}} = M_{\partial \mf{M}_{g,n}}. \]
如果考虑万有族 \(\mf{C}_{g,n}\)，我们可以构造 \(\pi^{-1}(\partial \mf{M}_{g,n}) \subseteq \mf{C}_{g,n}\)。因此，我们有映射
\[ \pi^* M_{\partial \mf{M}_{g,n}} \to M_{\pi^{-1}(\partial \mf{M}_{g,n})}. \]
我们也可以考虑除子 \(p_i(\mf{M}_{g,n}) \subset \mf{C}_{g,n}\)，它们同样定义对数结构
\[ P_i \coloneqq M_{p_i(\mf{M}_{g,n})}. \]
因此，我们可以定义对数结构
\[ M_{\mf{C}_{g,n}}^{\on{can}} \coloneqq M_{\pi^{-1}(\partial\mf{M}_{g,n})} \oplus_{\mc{O}^{\times}} \bigoplus_{i=1}^n P_i. \]
此时如果考虑
\[ \pi^{\flat} \colon \ul{\pi}^{-1} M_{\mf{M}_{g,n}}^{\on{can}} \to M_{\pi^{-1}(\partial \mf{M}_{g,n})} \to M_{\mf{C}_{g,n}}^{\on{can}}, \]
我们可以把 \(\pi\) 提升为对数栈之间的映射。

\begin{thm}[{\cite{FKato2000}}]
    映射 \(\pi\) 是对数光滑的整映射。
\end{thm}

\begin{defn}
    对于一个预稳定曲线族 \((\ul{C}/\ul{S}, \vec{p})\)，我们可以沿图表
    \begin{equation*}
    \begin{tikzcd}
        \ul{C} \ar{r} \ar{d} & \mf{C}_{g,n} \ar{d} \\
        S \ar{r} & \mf{M}_{g,n}
    \end{tikzcd}
    \end{equation*}
    拉回，从而给左边那条边赋予典范对数结构。
\end{defn}

\begin{rmk}
    设 \(\ul{S} = \Spec k\)。如果 \(S^{\on{can}} = \Spec (\N^r \to k)\)，那么 \(r\) 是 \(\ul{C}\) 上结点的个数。在一个结点 \(q\) 处平展局部地看，对数结构形如
    \begin{equation*}
    \begin{tikzcd}
        \N^{r-1} \oplus \N^2 \ar{r} & k[x,y]/xy \\
        \N^{r-1} \oplus \N \ar{u}{\on{id}\oplus\Delta} \ar{r} & k \ar{u}.
    \end{tikzcd}
    \end{equation*}
\end{rmk}

\begin{thm}[{\cite{FKato2000}}]
    设 \((\mc{C}/S, \vec{p})\) 是 fs 基底 \(S\) 上的 \(n\) 点标记对数曲线。则存在唯一的 \(M_S^{\on{can}} \to M_S\)，使得
    \begin{equation*}
    \begin{tikzcd}
        \mc{C} \ar{d} \ar{r} & \mc{C}^{\on{can}} \ar{d} \\
        S \ar{r}{a} & S^{\on{can}}
    \end{tikzcd}
    \end{equation*}
    是 fs 对数概形范畴中的笛卡儿图。
\end{thm}

一个口号是：对数曲线就是一条预稳定曲线再加上 \(M_S^{\on{can}} \to M_S\)。

\begin{exm}
    设 \(\ul{S} = \Spec k[s]\) 且 \(\ul{T} = \Spec k[t]\)。再取定 \(\ell \in \N\)。然后考虑图表
    \begin{equation*}
    \begin{tikzcd}
        \Spec \frac{k{ [x,y,s] }}{xy-s^{\ell}} \ar{d} \ar{r} & \Spec \frac{k{ [x,y,s] }}{xy-s^{\ell}} \ar{d} \ar{r} & \Spec  \frac{k{[x,y,t]}}{xy-t} \ar{r}\ar{d} & \mf{C}_{g,n} \ar{d} \\
        \Spec k{[s]} \ar{r}{\ul{ \on{id} }} & \Spec {k[s]} \ar{r} & \Spec k{[t]} \ar{r} & \mf{M}_{g,n}.
    \end{tikzcd}
    \end{equation*}
    在对数结构上，考虑左下箭头在 \(\N\) 上对应 \(1 \mapsto \ell\) 的情形。此时在左上角处，幺半群变为 \(\N^2 \oplus_{\N, \ell} \N\)。如果把它记作 \(P\)，那么映射 \(k[P] \to k[\N]\) 不是带典范对数结构的对数曲线族，但它是一个对数曲线族。
\end{exm}

\subsection{对数曲线的模}%
\label{sub:Moduli of log curves}

\begin{defn}
    \(n\) 点标记对数曲线之间的映射 \((\mc{C}, S) \to (\mc{C}', S')\) 是一对 \((\ell \colon \mc{C} \to \mc{C}', \Theta \colon S \to S')\)，使得 \(\Theta\) 是严格的且图表
    \begin{equation*}
    \begin{tikzcd}
        \mc{C} \ar{d} \ar{r}{\ell} & \mc{C}' \ar{d} \\
        S \ar{r}{\Theta} & S'
    \end{tikzcd}
    \end{equation*}
    是笛卡儿的。把参数化 \(n\) 点标记对数曲线的纤维范畴记作 \(\mf{M}_{g,n}^{\log}\)。
\end{defn}

\begin{prop}[{\cite{FKato2000,Olsson2003}}]
    我们有 \(\mf{M}_{g,n}^{\log} \simeq \ms{Log}_{(\mf{M}_{g,n}, M_{\mf{M}_{g,n}}^{\on{can}})}\)。
\end{prop}

特别地，对于一个概形 \(\ul{T}\)，群胚 \(\mf{M}_{g,n}^{\log}(\ul{T})\) 参数化对数结构 \(M_T\) 连同对数曲线 \(\mc{C} / T\)。

考虑 \(S = \Spec (Q \to k)\) 上的一条对数曲线 \(\pi \colon \mc{C} \to S, \vec{p}\)。记 \(\bar{\pi}^{\flat} \colon \bar{\pi}^* \bar{M}_S \to \bar{M}_{\mc{C}}\)。在泛点或一般光滑点 \(\eta\) 处，映射为
\[ \bar{\pi}_{\eta}^{\flat} \colon (\pi^{-1} \bar{M}_S)_{\eta} = Q \simeq \bar{M}_{\mc{C}, \eta}. \]
在标记点 \(p_i \in \ol{\ab\{\eta\}}\) 处，我们有 \(\bar{M}_{\mc{C}, p_i} = Q \oplus \N\)。此时我们有特化映射
\[ \chi_{\eta} \colon \bar{M}_{\mc{C}, p_i} \xrightarrow{\on{pr}_1} \bar{M}_{\mc{C}, \eta} \]
以及拉回
\[ \bar{\pi}_{p_i}^{\flat} \colon Q \hookrightarrow Q \oplus \N. \]
在结点 \(q\) 处，我们有 \(\bar{M}_{\mc{C}, q} = Q \oplus_{\N} \N^2\)。把 \(1\) 在 \(Q\) 中的像记作 \(\rho_q\)，它是 \(q\) 处的边长或光滑化参数。把 \(q\) 的两条分支的泛点记作 \(\bar{\eta}_1\) 与 \(\bar{\eta}_2\)。则泛化映射为
\begin{align*}
    \chi_{\eta_1} \colon Q \oplus_{\N} \N^2 \to Q \qquad (q, (a,b)) \mapsto q + a \rho_q ;\\
    \chi_{\eta_2} \colon Q \oplus_{\N} \N^2 \to Q \qquad (q, (a,b)) \mapsto q + b \rho_q .
\end{align*}

\subsection{对数曲线的热带化}%
\label{sub:Tropicalization of log curves}

我们将只考虑对数点 \(S = \Spec(Q \to k)\) 上的情形。我们将考虑锥的范畴 \(\ms{Cone}\)，其对象是满足 \(w \subseteq N_w \otimes_{\Z} \R\) 的对 \((w, N_w)\)，其中 \(w\) 是一个满维的强凸有理多面锥。其态射是映射 \((w_1, N_{w_1}) \to (w_2, N_{w_2})\)，由满足
\[ \chi \colon N_{w_1} \to N_{w_2} \]
使得 \(\chi_{\R}(w_1) \subseteq w_2\) 的映射组成。

\begin{exm}
    设 \(X\) 为 fs 对数概形，考虑 \(( (\bar{M}_{X,x}^{\vee})_{\R}, (\bar{M}_{X,x}^{\vee})_{\Z} )\)。
\end{exm}

现在我们将考虑广义锥复形的范畴 \(\ms{GCC}\)。其对象是拓扑空间 \(\abs{\Sigma}\) 连同把它表示为 \(\ms{Cone}\) 中一个所有箭头都是面的包含的图表之余极限的展示。其态射是连续映射 \(\abs{\Sigma} \to \abs{\Sigma'}\)，使得对每个 \((w, N_w) \in \Sigma\)，映射 \(w \to \abs{\Sigma'}\) 经由某个 \((w, N_w) \to (w', N_{w'})\) 分解。

\begin{exm}
    如果 \(f \colon X \to Y\) 是 fs 对数概形的态射，则映射 \(\bar{f}^{\flat} \colon f^* \bar{M}_Y \to \bar{M}_X\) 在热带化上诱导 \(\Sigma(X) \to \Sigma(Y)\)。
\end{exm}

\begin{defn}
    \textit{锥复形}是一个广义锥复形 \(\Sigma\)，使得对每个 \((\sigma, N_{\sigma}) \in \Sigma\)，映射 \(\sigma \to \abs{\Sigma}\) 是单射。
\end{defn}

\begin{thm}[{\cite{ACP2015}}]
    如果 \(X\) 是 fs 的、在 \(\ul{\C}\) 上对数光滑且带有 Zariski 对数结构，那么 \(\Sigma(X)\) 是一个锥复形。
\end{thm}

对于一个幺半群 \(Q\)，我们将记 \(w_Q = Q_{\R}^{\vee} = \Hom(Q, \R_{\geq 0})\)，并记格 \(\Hom(Q, \Z) = Q_{\Z}^{\vee}\)。在光滑点处，我们有
\[ w_{\eta} = (Q_{\R}^{\vee}, Q_{\Z}^{\vee}). \]
在标记点处，我们看到
\[ w_p = \ab(\Hom (Q \oplus \N, \R_{\geq 0}) = Q_{\R}^{\vee} \oplus \R_{\geq 0}, Q_{\Z}^{\vee}\oplus \Z). \]
在结点处，我们有
\begin{align*}
    w_q &= \ab(\Hom(Q \oplus_{\N} \N^2, \R_{\geq 0}), ( \bar{M}_{\mc{C}, q} )_{\Z}^{\vee}).
\end{align*}
注意这个锥实际上等于
\[ \ab\{(m, m_x, m_y) \mid m(\rho_q) = m_x + m_y \} \simeq \ab\{(m, e_x) \mid e_x \leq m(\rho_q) \} \subseteq Q_{\R}^{\vee} \times \R_{\geq 0}. \]

既然我们已经了解了与每个点相伴的锥，现在需要讨论面映射。对应于标记点的泛化映射是
\[ Q \oplus \N \xrightarrow{\on{pr}_1} Q, \]
相应的面映射是
\[ w_{\eta} \hookrightarrow w_p \qquad m \mapsto (m,0). \]
在结点处，两个面映射是
\[ \chi_{\eta_1}^{\vee} \colon Q_{\R}^{\vee} \to w_q \qquad m \mapsto (m, m(\rho_q)) \]
与
\[ \chi_{\eta_2}^{\vee} \colon Q_{\R}^{\vee} \to w_q \qquad m \mapsto (m, 0). \]

\begin{exm}
    考虑一条曲线，它是两条有理曲线之并，其中一条分支上有一个标记点 \(p\)。令 \(\ul{S} = \Spec \C\)。则我们看到
    \[ S^{\on{can}} = \Spec(\N \to \C), \]
    \(Q = \N\)，且 \(Q_{\R}^{\vee} = \R_{\geq 0}\)。我们还看到 \(\bar{M}_{C, q} = \N \oplus_{\N} \N^2\) 且 \(\bar{M}_{C, p} = \N \oplus \N\)。
\end{exm}

\begin{figure}[H]
    \centering
    \begin{tikzpicture}[
        x=1cm,
        y=1cm,
        cone ray/.style={thick,->,>=stealth},
        common face/.style={red,line width=2pt,line cap=round},
        diagonal face/.style={blue!75!black,line width=2pt,line cap=round}
    ]
        % The cone at the marked point.
        \coordinate (p) at (0,4.45);
        \draw[cone ray] (p) -- ++(-0.55,2.15);
        \draw[cone ray] (p) -- ++(2.25,1.85);
        \draw[diagonal face]
            ($(p)+(0.75,0.62)$) -- ($(p)+(1.45,1.19)$);
        \node[above] at ($(p)+(0.55,2.2)$)
            {$w_p=\R_{\geq 0}\times\R_{\geq 0}$};

        % The cone at the first generic point.
        \coordinate (etaone) at (0.35,4);
        \draw[cone ray] (etaone) -- ++(2.25,1.85);
        \draw[diagonal face]
            ($(etaone)+(0.75,0.62)$) -- ($(etaone)+(1.45,1.19)$);
        \node[above right] at ($(etaone)+(2.0,1.6)$) {$w_{\eta_1}$};

        % The cone at the node.
        \coordinate (q) at (0.35,3.25);
        \draw[cone ray] (q) -- ++(2.25,1.85);
        \draw[cone ray] (q) -- ++(3.35,0);
        \draw[diagonal face]
            ($(q)+(0.75,0.62)$) -- ($(q)+(1.45,1.19)$);
        \draw[common face] ($(q)+(0.85,0)$) -- ($(q)+(1.75,0)$);
        \node at ($(q)+(2.55,0.75)$) {$w_q$};
        \node[anchor=west] at ($(q)+(4.05,0.75)$)
            {$\bigl\{(a,b)\mid b\leq a\bigr\}\subseteq\R_{\geq 0}^{2}$};

        % The colored intervals indicate the common faces used in the gluing.
        \coordinate (etatwo) at (0.35,2.75);
        \draw[cone ray] (etatwo) -- ++(3.35,0);
        \draw[common face] ($(etatwo)+(0.85,0)$) -- ($(etatwo)+(1.75,0)$);
        \node[right] at ($(etatwo)+(3.45,0)$) {$w_{\eta_2}$};
    \end{tikzpicture}
\end{figure}


\subsection{热带曲线}%
\label{sub:Tropical curves}

\begin{defn}
    锥 \(w, N_w\) 上的一族\textit{热带曲线} \((G, \vec{g}, \vec{\ell})\) 由以下组成：
    \begin{enumerate}
        \item 一个连通图 \(G = V(G) \cup E(G) \cup L(G)\)；
        \item 一个腿序 \(\vec{m} \colon L(G) \xrightarrow{\sim} \ab\{1, \ldots, n\}\)；
        \item 一个亏格函数 \(\vec{g} \colon V(G) \to \N\)；
        \item 一个边长函数 \(\ell \colon E(G) \to \Hom(w \cap N_w, \N) \setminus \ab\{0\}\)。
    \end{enumerate}
\end{defn}

现在我们将构造 \(\Sigma(G, \ell)\)。对 \(v \in V(G)\)，令 \(w_v = w\)。对 \(e \in E(G)\)，令
\[ w_e = \ab\{ (s, \lambda) \mid \lambda \leq \vec{\ell}(e)(s) \} \subseteq w \times \R_{\geq 0}. \]
对腿 \(\ell \in L(G)\)，令 \(w_{\ell} = w \times \R_{\geq 0}\)。现在我们需要定义面映射。如果 \(\ell\) 与 \(v\) 关联，则面映射为
\[ w_v \hookrightarrow w_{\ell} \colon s \mapsto (s, 0). \]
如果 \(e\) 与两个顶点 \(v_1, v_2\) 关联，则面映射为
\begin{align*}
    w_{v_1} \hookrightarrow w_e &\colon s \mapsto (s, \vec{\ell}(e)(s)) \\
    w_{v_2} \hookrightarrow w_e & \colon s \mapsto (s, 0).
\end{align*}
\begin{rmk}
    \(v_1\) 与 \(v_2\) 的选取在同构意义下不改变锥复形，但当我们要为对数映射定义接触阶时，它确实有影响。然而一般来说，没有典范的方式作出这一选择。
\end{rmk}

\subsection{穿孔曲线}%
\label{sub:Punctured curves}

穿孔对数曲线与映射在~\cite{ACGS2025} 中建立。

我们考虑结点曲线
\[ \Spec k[x,y]/xy \to \Spec k \]
其对数结构由带图卡 \((a,b) \mapsto x^a y^b\) 的 \(\N^2\) 给出，而 \(\N \to k\) 为标准对数点。现在考虑 \(\ul{C}_1 = Z(x) = \Spec k[y]\)，并用 \(p \colon \Spec k \to \Spec k[y]\) 标出该结点。注意
\[ C^{\on{can}} = (\ul{C}, M_C) \to S^{\on{can}} = \Spec(\N \to k). \]
记 \(M_{C_1}^{\circ} = M_C|_{C_1}\)。于是我们看到在对数结构上有图表
\begin{equation*}
\begin{tikzcd}
    k[y] & & \N^2 \ar{ll}[swap]{(a,b) \mapsto 0^a y^b} \\
    k \ar{u} & & \N \ar{ll} \ar{u}{\Delta}.
\end{tikzcd}
\end{equation*}
特别地，我们可以计算茎 \(\bar{M}_{C_1}^{\circ}|_p = \N^2\)，因此我们有
\[ \bar{\pi}_{C_1^{\circ}}^{\flat} |_p \colon \N \xrightarrow{\Delta} \N^2. \]
然而，我们本应有映射
\[ \bar{\pi}^{\flat}_{C_1}|_p \colon \N \xrightarrow{i_1} \N \oplus \N, \]
因此 \(C_1\)\textbf{不是}对数曲线。于是我们可以考虑图表
\begin{equation*}
\begin{tikzcd}
    \N \oplus \Z \\
    \N^2 \ar[hookrightarrow]{u} \ar{rr}{(0,1)\mapsto (0,1)}[swap]{(1,0)\mapsto (1,1)} & & \N^2 \ar{ull}[swap]{\Vectorstack{{(1,0)\mapsto(1,-1)} {(0,1)\mapsto(0,1)}}} \\
    & \N \ar{ul}{i_1} \ar{ur}{\Delta}.
\end{tikzcd}
\end{equation*}
因此，我们将考虑扩张对数结构。

回忆对于对数曲线，我们有 \(M_{\mc{C}} = M_{\mc{C}}' \oplus_{\mc{O}_{\mc{C}}^{\times}} P_{\mc{C}}\)，其中 \(P_{\mc{C}}\) 是由标记给出的除子对数结构。

\begin{defn}
    fs 对数概形 \(S\) 上的\textit{穿孔曲线}由 \((\mc{C}^{\circ} \xrightarrow{p} \mc{C} \xrightarrow{\pi} S, \{p_i\}_{i=1}^n)\) 组成，其中 \((\mc{C} \xrightarrow{\pi} S, \vec{p})\) 是一条对数曲线，而映射 \(p \colon \mc{C}^{\circ} \to \mc{C}\) 是细对数概形的映射，满足以下条件：
    \begin{enumerate}
        \item 底层映射 \(\ul{p}\) 是同构；
        \item 在对数结构上，映射 \(p^{\flat} \colon M_{\mc{C}} \hookrightarrow M_{\mc{C}^{\circ}}\) 在 \(p_i\) 之外是同构，并诱导
            \[ M_{\mc{C}} \xrightarrow{p^{\flat}} M_{\mc{C}^{\circ}} \hookrightarrow M_{\mc{C}} \oplus P^{\mr{gp}}. \]
        \item 对任意标记 \(p_i\) 与任意 \(s \in M_{\mc{C}^{\circ}, p_i} \setminus M_{\mc{C}, p_i}\)，我们有 \(\alpha_{\mc{C}^{\circ}}(s) = 0\)。
    \end{enumerate}
\end{defn}

粗略地说，穿孔曲线就是一条对数曲线再加上一个交换图表
\begin{equation*}
\begin{tikzcd}
    M_{\mc{C}^{\circ}} \ar{dr} \ar[hookrightarrow]{r} & M_{\mc{C}}' \oplus_{\mc{O}_{\mc{C}}^{\times}} P^{\on{gp}} \\
    & \mc{O}_{\mc{C}} \\
    M_{\mc{C}} \ar[hookrightarrow]{uu} \ar{ur}.
\end{tikzcd}
\end{equation*}
的数据。


\begin{exm}
    设 \((\ul{\mc{C}} \to \ul{S} = \Spec k, p)\) 为一条光滑的 \(1\) 点标记曲线。
    \begin{enumerate}
        \item 令 \((\mc{C}/S, p)\) 带有典范对数结构（于是 \(S\) 没有对数结构）。设 \(M_{\mc{C}} = P\)。则任何穿孔 \(p^{\flat} \colon M_{\mc{C}} \hookrightarrow M_{\mc{C}^{\circ}} \subset P^{\on{gp}}\) 必为平凡的。原因是 \(M_{\mc{C}^{\circ}}\) 不可能有满足 \(\alpha_{\mc{C}^{\circ}}(s) = 0\) 的截面 \(s\)。如果 \(s \in P = M_{\mc{C}}\)，则 \(\alpha_{\mc{C}^{\circ}}(s) = \alpha_{\mc{C}}(s)\)。否则，如果 \(s \notin P\)，那么由 \(\alpha_{\mc{C}^{\circ}}(-s)\) 非零可知 \(\alpha_{\mc{C}^{\circ}}(s)\) 也非零。
        \item 设 \(M_{\mc{C}^{\circ}} \subseteq (\N \oplus \mc{O}_{\mc{C}}^{\times}) \oplus P^{\on{gp}}\) 是一个穿孔。取 \(\mc{O}(-p)\) 的一个局部生成元 \(x\)。则任意 \(s \in M_{\mc{C}^{\circ}, p}\) 都具有形式 \(s = ((n, \varphi), x^m)\)，其中 \(n \in \N\) 且 \(\varphi \in \mc{O}_{\mc{C}, p}^{\times}\)。如果 \(m < 0\)，那么我们要求 \(\alpha_{\mc{C}^{\circ}}(s) = 0\)。这意味着 \(\alpha_{\mc{C}}((n, \varphi), 0) = 0\)，因此
            \[ \bar{M}_{\mc{C}^{\circ}, p} \subseteq \ab\{(n, m) \in \N \oplus \Z \mid m \geq 0 \text{ if } n = 0 \}. \]
            事实上，右端任何包含 \(\N \oplus \N\) 的细子幺半群都可以由某个穿孔实现。
        \item 令 \(S = \Spec(\N \to k)\)。考虑满足 \(\bar{M}_{\mc{C}^{\circ}, p} = \N^2 + \N (2,-2)\) 的 \(M_{\mc{C}^{\circ}}\)。由于 \((1,-1) \notin \bar{M}_{\mc{C}^{\circ}, p}\)，我们看到 \(M_{\mc{C}^{\circ}}\) 不是饱和的。
    \end{enumerate}
\end{exm}

\begin{lem}[{\cite{ACGS2025}}]
    设 \(Q\) 为尖锐的 fs 幺半群。记 \(w \coloneqq \Hom(Q, \R_{\geq 0}) = Q_{\R}^{\vee}\)。设 \(Q^{\circ} \subseteq Q \oplus \Z\) 有限生成，满足 \(Q \oplus \N \subsetneq Q^{\circ}\) 且 \(Q^{\circ} \cap \ab\{0\} \times \Z_{<0} = \emptyset\)。则存在一个非零的、凹的、斜率为有理数的分段线性函数
    \[ \ell \colon w \to \R_{\geq 0} \]
    使得
    \[ (Q^{\circ})_{\R}^{\vee} = \ab\{ (s, \lambda) \in w \times \R_{\geq 0} \mid 0 \leq \lambda \leq \ell(s) \}. \]
\end{lem}

\begin{exm}
    令 \(Q = \N^2\)，并设
    \begin{align*}
        Q^{\circ} &= \ab<((1,0), -1), ((0,1), -1), ((0,0),1)> \subset \N^2 \oplus \Z \\
        &= \ab\{((a,b), m) \mid a+b+m \geq 0, a,b \geq 0 \}.
    \end{align*}
    则
    \[ (Q^{\circ})_{\R}^{\vee} = \ab\{ (e_x, e_y, e_t) \mid 0 \leq e_t \leq \min(e_x, e_y) \}. \]
\end{exm}
